% Auxiliary bounds, not a finite-generation claim or a new requested-row closure.
\begin{proposition}[Proved: arithmetic genus of the full exceptional fibre]
\label{D:exceptional-arithmetic-genus-bound}
Use the one-centre setup of
Propositions~\ref{D:surface-uniform-different} and
\ref{D:fixed-boundary-uniform-different}. Thus $Y\to(S,s)$ is a
fixed smooth projective resolution of a normal surface germ over
an algebraically closed field, $A=\sum_i a_iE_i$ is a fixed
positive exceptional cycle with $A\cdot E_i<0$, and $D$ is the
strict transform of the original reduced divisorial boundary.
Set $D=0$ in the punctured case.

Let a connected normal finite covering in the preceding
etale or tame setup have degree $n$. In particular its extension
of function fields is separable. Resolve its normalization over
$Y$ and contract only the smooth rational $(-1)$-curves mapped
to points of $Y$, as in the cited propositions. Write
\[
 f:Z\longrightarrow Y,\qquad C=f^*A,\qquad
 F=C_{\mathrm{red}},\qquad c=\#\pi_0(F),\qquad
 b=(K_Y+D)\cdot A+\sum_i a_i.
\]
Then $c\le n$ and
\[
             2h^1(F,\mathcal O_F)\le nb+c\le n(b+1).    \tag{D.AG.1}
\]
The possibly nonreduced proper Cartier divisor $C$ satisfies
\[
 2+n\left(A^2-\sum_i a_i\right)
 \le 2p_a(C)
 \le 2+n\left(A^2+(K_Y+D)\cdot A\right).               \tag{D.AG.2}
\]
If we further resolve the reduced exceptional fibre by point
blowups to obtain a strict normal crossings fibre $\widehat F$,
then
\[
 \sum_{G\subset\widehat F}g(G)+b_1(\Gamma_{\widehat F})
                   \le \frac{nb+c}{2},                 \tag{D.AG.3}
\]
where $\Gamma_{\widehat F}$ is its dual multigraph.
\end{proposition}
\begin{proof}
All curves in the support of $C$ are exceptional over the finite
set of source points above $s$. Conversely every such exceptional
curve belongs to this support, since $A$ is positive on all
components over $s$. Write $F_{i\lambda}$ for the curves
dominating $E_i$, and $Q_j$ for those mapped to points of $Y$.
The coefficient of $C$ along $F_{i\lambda}$ is
$a_i e_{i\lambda}$. Its coefficients along the $Q_j$ are
positive integers.

The construction and adjunction estimates in the proof of
Proposition~\ref{D:surface-uniform-different} give
\[
 K_Z\cdot Q_j\ge0,\qquad K_Z\cdot F_{i\lambda}\ge-1,
 \qquad K_Z\cdot C\le n(K_Y+D)\cdot A.                 \tag{D.AG.4}
\]
The last inequality uses the effective determinant-of-differentials
divisor. Its components over the $E_i$ have nonpositive
intersection with $C$, its point-mapped components have zero
intersection, and its horizontal part has intersection at most
$nD\cdot A$ by original-boundary tameness. This argument does
not require the cover to be Galois.

In passing from $C$ to its reduced divisor, lowering a coefficient
on a point-mapped component cannot increase the canonical
intersection. Lowering the coefficient $a_i e_{i\lambda}$ on a
dominating component increases it by at most
$a_i e_{i\lambda}-1$. Therefore
\[
 \begin{split}
 K_Z\cdot F
 &\le K_Z\cdot C+
                \sum_{i,\lambda}(a_i e_{i\lambda}-1)\\
 &\le n(K_Y+D)\cdot A+n\sum_i a_i=nb,
 \end{split}                                          \tag{D.AG.5}
\]
where $\sum_\lambda e_{i\lambda}\le n$ follows from
$\sum_\lambda e_{i\lambda}f_{i\lambda}=n$ with full residue
degrees.

The exceptional intersection matrix is negative definite
\citep[Theorem~10.1]{Kollar2013D}. Distinct connected components
of $F$ are disjoint. Each has a negative integral square,
so $F^2\le-c$. Since $F$ is proper and reduced over the
algebraically closed field, $h^0(\mathcal O_F)=c$.
Adjunction for a proper Cartier divisor on the smooth surface
therefore gives
\[
 2h^1(\mathcal O_F)
   =2c+F^2+K_Z\cdot F\le nb+c.
\]
To check $c\le n$, pass to the strict henselization at $s$.
The finite normalization has at most $n$ local normal factors:
their positive generic degrees add to $n$. The exceptional
fibre of the proper birational morphism over each such normal
local factor is connected, by its pushforward of the structure
sheaf and Stein factorization. This proves the component bound.

For (D.AG.2), projection gives $C^2=nA^2$. The lower canonical
bound from the same construction is
$K_Z\cdot C\ge-n\sum_i a_i$. Combining this and (D.AG.4) with
$2p_a(C)=2+C^2+K_Z\cdot C$ proves both inequalities.

Finally consider one blowup of a point of a reduced divisor of
multiplicity $m\ge1$. Its full reduced total transform is
\[
                 F'=\pi^*F-(m-1)E,
\]
where $E$ is the exceptional $(-1)$-curve. The equalities
$K_{Z'}=\pi^*K_Z+E$, $E^2=-1$ and
$E\cdot\pi^*F=0$ give
\[
             p_a(F')=p_a(F)-\frac{(m-1)(m-2)}2.
\]
The number of connected components is unchanged. Hence $h^1$
of the reduced full fibre does not increase under these blowups.
For a reduced strict normal crossings curve, its normalization
exact sequence gives
$h^1(\mathcal O_{\widehat F})=\sum_G g(G)+b_1(\Gamma_{\widehat F})$.
This proves (D.AG.3). Covers of the punctured strict henselization
are treated by the same finite-presentation descent used in the
cited different bounds.
\end{proof}

\begin{lemma}[Proved: Frobenius saturation of a finite normalization lattice]
\label{D:normalization-Frobenius-saturation}
Let $A$ be a complete discrete valuation ring of characteristic
$p>0$, with uniformizer $\pi$ and fraction field $K$. Let $B$
be its integral closure in a finite separable extension $L/K$.
Write $e$ for the ramification index and $\delta$ for the
different exponent. For every power $q$ of $p$,
\[
  \pi^{\lceil(q-1)\delta/e\rceil}B
                     \subset A\cdot B^q\subset B.      \tag{D.AG.6}
\]
No separability of the residue extension is assumed.

Consequently, let $\rho:Y'\to Y$ be finite, dominant and
generically separable between normal surfaces in characteristic
$p$, with $Y$ smooth. Put $\mathcal B=\rho_*\mathcal O_{Y'}$.
Suppose all ramification is supported on fixed prime divisors
$D_j$, and at every prolongation of $D_j$ we have
$\delta/e\le c_j$. Then $\mathcal B$ is locally free, and the
relative Frobenius map gives actual inclusions
\[
 \mathcal B(-H)\subset
       \operatorname{im}(F_Y^*\mathcal B\longrightarrow\mathcal B)
       \subset\mathcal B,
 \qquad H=\sum_j\lceil(p-1)c_j\rceil D_j.              \tag{D.AG.7}
\]
The divisor $H$ is independent of the degree of the covering.
\end{lemma}
\begin{proof}
Choose an $A$-basis $b_1,\ldots,b_n$ of $B$ and its trace-dual
$K$-basis $b_1^*,\ldots,b_n^*$ of $L$. The inverse different is
the trace-dual lattice, so $v_B(b_i^*)\ge-\delta$; this is the
trace construction of the different
\citep[Section~49.8]{StacksDDifferent}. Since
$\operatorname{Tr}(x^q)=\operatorname{Tr}(x)^q$ in characteristic
$p$, the bases $b_i^q$ and $(b_i^*)^q$ are also trace-dual.
For $b\in B$ we therefore have
\[
        b=\sum_i\operatorname{Tr}_{L/K}
                         \bigl(b(b_i^*)^q\bigr)b_i^q.
\]
Set $M=\lceil(q-1)\delta/e\rceil$. The element
$\pi^M b(b_i^*)^{q-1}$ is integral, because its valuation
is at least $eM-(q-1)\delta\ge0$. Pairing it with $b_i^*$
therefore has integral trace. Thus
\[
 \pi^M\operatorname{Tr}_{L/K}\bigl(b(b_i^*)^q\bigr)\in A
\]
for every $i$, proving (D.AG.6). The proof works componentwise
if the completed normalization is a product of DVRs, with
$M$ the largest corresponding ceiling.

A normal surface is Cohen--Macaulay. Miracle flatness over the
regular surface $Y$ therefore makes $\mathcal B$ locally free
\citep[Tag 00R4]{StacksArithmeticDefectB}. The map
$F_Y^*\mathcal B\to\mathcal B$, $a\otimes b\mapsto ab^p$,
is generically an isomorphism, by separability of the function
field extension. At every codimension-one point, completion
and (D.AG.6) show that its inverse carries
$\mathcal B(-H)$ into $F_Y^*\mathcal B$.
Both are vector bundles on the normal scheme $Y$. A rational
morphism between them which is regular at every codimension-one
point extends everywhere, by the intersection description of
reflexive sheaves \citep[Tag 0EBJ]{C_StacksPicNorm}. This gives the first actual sheaf inclusion
in (D.AG.7); the second is the original Frobenius map. In
particular we have not replaced its image by a reflexive hull
without justification.
\end{proof}

\begin{remark}[Limits of these two bounds]
The arithmetic bound retains dual-graph cycles and normalized
genera but permits arbitrarily long rational trees. Contracting
such trees can create normalization defects on curves in a
finite normal model, so (D.AG.1) does not remove the nilradical
term in the nonreduced-fibre calculation following
Lemma~\ref{D:proper-curve-arithmetic-defect}.
The Frobenius saturation bound concerns the actual canonical
normalization lattice. It does not by itself give uniform
Harder--Narasimhan windows on the exceptional components,
an atlas with size linear in the degree of the cover, or a
factorial bound for the number of finite covers. None of these
statements asserts finite generation of the full group.
\end{remark}
